\chapter{研究设计}

\section{资料来源与纳入原则}

本示例完全基于公开资料，不包含真实病历、内部访谈录音或未公开论文。资料来源包括媒体报道、节目档期、公号运营记录、杂志封面时间线以及公开可检索的互动数据。纳入标准强调事件时间明确、信息来源可追溯，且能够对应到职业阶段变化。

\section{分析框架}

研究采用“阶段划分-事件编码-韧性分析”的三层框架。阶段划分解决“何时发生转折”，事件编码回答“转折由什么驱动”，韧性分析则讨论危机发生后，主体如何通过平台迁移与叙事调整重新建立稳定表达。

\section{示例项目的工程意义}

对模板仓库而言，公开案例研究的意义在于：既能展示摘要、目录、缩略词、章节和参考文献的完整写作结构，又能在不接触任何保密正文的前提下，让样式验证和构建回归长期保留在公开仓库中。
